% !PRINTED 363
\weberpage{363}

%<<CONT
man die Ideale nach Bd.~II, \S\,169 durch die in Bd.~II, \S\,153\,ff. besprochenen \gesperrt{Funktionale} repräsentiert, weil man dabei die gewöhnlichen Regeln der Multiplikation und Division benutzen kann\footnote{Vgl. auch des Verfassers Abhandlung \glqq Über Zahlengruppen in algebraischen Körpern\grqq{}, erste Mitteilung, Mathematische Annalen 48, 438.}.

\gesperrt{Wir beschränken uns hier ein für allemal in dem Körper $\Omega$ auf Zahlen und Ideale, die zu einer beliebig anzunehmenden ganzen rationalen Zahl $Q$} (dem Führer einer Ordnung) \gesperrt{teilerfremd sind.}

Nach dieser Voraussetzung bilden die ganzen und gebrochenen Zahlen des Körpers $\Omega$ die vorher mit $O$ bezeichnete Gruppe. Wir erweitern diese Gruppe durch Hinzufügung der \gesperrt{Funktionale} und bezeichnen die so erweiterte \gesperrt{Gruppe} mit $\overline{O}$. Jedem Element von $\overline{O}$ entspricht dann ein ganzes oder gebrochenes Ideal, und solchen Elementen, deren Quotient eine funktionale Einheit ist, \gesperrt{entspricht dasselbe Ideal.}

Ist $\eta$ eine Zahl in $O$ und $\varepsilon$ irgend eine funktionale Einheit, so entspricht dem Produkt $\varepsilon\eta$ ein \gesperrt{Hauptideal}. Demnach bezeichnen wir
\begin{quote}
\gesperrt{mit $\overline{E}$ die Gruppe der funktionalen Einheiten,}\\
\gesperrt{\glqq\ $E$\ \glqq\ \glqq\ \glqq\ numerischen Einheiten,}
\end{quote}
und durch
\begin{quote}
\gesperrt{$\overline{E}O$ die Hauptklasse der Funktionale.}
\end{quote}
\gesperrt{Nehmen wir ein Repräsentantensystem}
\weq{1}{\varphi_1,\ \varphi_2,\ \dots,\ \varphi_h}
nicht äquivalenter ganzer Funktionale in $\Omega$, so wird jedes ganze oder gebrochene Funktional $\Phi$ in der Weise darstellbar sein:
\[
\Phi = \varphi_i \overline{\omega},
\]
wo $\overline{\omega}$ ein Funktional aus $\overline{E}O$ ist, und die Idealklassen in $\Omega$ sind die Nebengruppen von $\overline{E}O$ in $\overline{O}$. Demnach ist die Klassenzahl $h$ der Index des Teilers $\overline{E}O$ von $\overline{O}$:
\weq{2}{h = (\overline{O},\ EO).} % [sic] Gedruckt ist E hier ohne Überstrich.
Hierbei hat man sich bei der schärferen Klasseneinteilung bei $O$ auf die Zahlen mit positiver Norm zu beschränken.

Betrachten wir nun die Klasseneinteilung der Ideale nach der Ordnung $[Q]$, so haben wir als \gesperrt{Hauptklasse} die Funktionalgruppe $\overline{E}O'$ zu betrachten und die Klassenzahl ist
\[
h' = (\overline{O},\ \overline{E}O').
\]
